% !TeX root = ../analysis1.tex
\chapter{Foundations and Sequences}\label{chap:analysis-foundations}

\section{What analysis is about}
Analysis provides precise definitions and proofs for ideas from calculus:
limits, infinite sums, continuity and derivatives. An \emph{axiom} is a
starting assumption; a \emph{theorem} is a statement proved from assumptions.
A \emph{lemma} is a result used to establish another result, and a
\emph{corollary} follows readily from a previous result.

\begin{example}[Why infinite sums need definitions]
The harmonic series $\sum_{n=1}^{\infty}1/n$ diverges, whereas
$\sum_{n=1}^{\infty}(-1)^{n+1}/n$ converges. Rearranging a conditionally
convergent series may change its sum. Thus the rules for finite sums cannot
be applied to infinite sums without checking their hypotheses.
Chapter~\ref{chap:analysis-series} makes these statements precise.
\end{example}

\begin{example}[Smooth does not imply equal to a Taylor series]
Consider
\[
f(x)=\begin{cases}e^{-1/x^2},&x>0,\\0,&x\leq0.\end{cases}
\]
All derivatives at $0$ exist and are zero, so its Taylor series there is the
zero series. Nevertheless $f(x)>0$ for $x>0$. Merely having derivatives of
all orders does not ensure that a function equals its Taylor series.
This is a motivating example; the differentiability assertion is revisited
at the end of Chapter~\ref{chap:analysis-differentiation}.
\end{example}

\section{Number systems and rational approximation}
\subsection{From natural numbers to real numbers}
The familiar number systems satisfy
\[
\mathbb N\subset\mathbb Z\subset\mathbb Q\subset\mathbb R\subset\mathbb C.
\]
Here $0\notin\mathbb N$. Addition and multiplication on $\mathbb N$ are
commutative and associative, and multiplication distributes over addition:
\begin{align*}
m+n&=n+m,&mn&=nm,\\
(m+n)+p&=m+(n+p),&(mn)p&=m(np),\\
m(n+p)&=mn+mp.
\end{align*}
For any $m,n$, exactly one of $m<n$, $m=n$, $m>n$ holds. Order is transitive.

Subtraction need not stay in $\mathbb N$: $x+m=n$ has the integer solution
$x=n-m$. This motivates
$\mathbb Z=\{\ldots,-2,-1,0,1,2,\ldots\}$.
Division need not stay in $\mathbb Z$, motivating
\[
\mathbb Q=\left\{\frac pq:p,q\in\mathbb Z,\ q\ne0\right\}.
\]
Fractions have the rules
\[
\frac mn+\frac pq=\frac{mq+pn}{nq},\qquad
\frac mn\frac pq=\frac{mp}{nq},\qquad
\frac{m/n}{p/q}=\frac{mq}{np}\quad(p\ne0).
\]
Every rational has a reduced representation with coprime numerator and
denominator. Terminating or eventually periodic decimal expansions represent
rational numbers. Non-periodic decimal expansions motivate a larger system,
$\mathbb R$; a real number outside $\mathbb Q$ is called \emph{irrational}.
We will use completeness, rather than decimals, to characterise $\mathbb R$.

\begin{lemma}[Parity]
If an integer $n$ is odd, then $n^2$ is odd. Consequently, if $n^2$ is even,
then $n$ is even.
\end{lemma}
\begin{proof}
Writing $n=2k+1$ gives
$n^2=4k^2+4k+1=2(2k^2+2k)+1$. The second statement is the contrapositive.
\end{proof}

\begin{theorem}[Irrationality of $\sqrt2$]
There is no rational number $x$ with $x^2=2$.
\end{theorem}
\begin{proof}
Suppose $x=p/q$ in lowest terms, with $q>0$. Then $p^2=2q^2$, so $p$ is even.
Write $p=2k$. Substitution gives $4k^2=2q^2$, hence $q^2=2k^2$, and $q$ is
also even. This contradicts coprimality.
\end{proof}

\subsection{Pell's equation and approximation}
Once the positive square root of $2$ is available, the equation
$|p^2-2q^2|=1$, for $p,q\in\mathbb N$, gives
\[
\left|\frac pq-\sqrt2\right|
=\frac{1}{q^2(p/q+\sqrt2)}<\frac1{q^2}.
\]
Thus a solution with large $q$ gives a good rational approximation.

\begin{proposition}[A recursive family of approximations]
Define $p_1=q_1=1$ and
\[
p_{n+1}=p_n+2q_n,\qquad q_{n+1}=p_n+q_n.
\]
Then $p_n,q_n\in\mathbb N$, $p_n^2-2q_n^2=(-1)^n$, and $q_n\geq n$.
Consequently $p_n/q_n\to\sqrt2$.
\end{proposition}
\begin{proof}
The assertions hold at $n=1$. If they hold at $n$, then
\begin{align*}
p_{n+1}^2-2q_{n+1}^2
&=(p_n+2q_n)^2-2(p_n+q_n)^2\\
&=-(p_n^2-2q_n^2)=(-1)^{n+1},
\end{align*}
and $q_{n+1}\geq1+n$. Induction proves the claims. The error is less than
$1/q_n^2\leq1/n^2$, which tends to zero.
\end{proof}
\begin{center}
\begin{tabular}{rrrr}
\toprule
$n$&$p_n$&$q_n$&$p_n/q_n$ (approximately)\\\midrule
1&1&1&1.000000\\
2&3&2&1.500000\\
3&7&5&1.400000\\
4&17&12&1.416667\\
5&41&29&1.413793\\
6&99&70&1.414286\\\bottomrule
\end{tabular}
\end{center}
\begin{remark}
An induction proof checks a base case and proves that the assertion for $n$
implies the assertion for $n+1$. The limiting conclusion above completes an
argument left unfinished on source p.~9.
\end{remark}

\section{Ordered fields and inequalities}
\begin{definition}[Field]
A field $F$ is a set with addition and multiplication satisfying the following
properties for all $a,b,c\in F$:
\begin{enumerate}
\item $(a+b)+c=a+(b+c)$ and $a+b=b+a$.
\item There is $0\in F$ with $a+0=a$; every $a$ has an additive inverse $-a$.
\item $(ab)c=a(bc)$ and $ab=ba$.
\item There is $1\in F$, $1\ne0$, with $a1=a$; every nonzero $a$ has an
inverse $a^{-1}$ with $aa^{-1}=1$.
\item $a(b+c)=ab+ac$.
\end{enumerate}
These are the nine arithmetic axioms, grouped by operation.
\end{definition}

\begin{definition}[Ordered field]
An ordered field is a field with a total order $\leq$ satisfying
antisymmetry and transitivity, together with
\[
a\leq b\ \Longrightarrow\ a+c\leq b+c,
\qquad
a\leq b,\ c\geq0\ \Longrightarrow\ ac\leq bc.
\]
Totality means that $a\leq b$ or $b\leq a$ for every pair. Multiplication by
a negative number reverses the inequality.
\end{definition}
Both $\mathbb Q$ and $\mathbb R$ are ordered fields. The condition $c\geq0$
is essential; it was omitted from the corresponding axiom in the source.

\begin{lemma}[Basic field consequences]
In a field, $a0=0$. Also, $ab=0$ implies $a=0$ or $b=0$.
\end{lemma}
\begin{proof}
Distributivity gives $a0=a(0+0)=a0+a0$; subtract $a0$. If $a\ne0$, multiply
$ab=0$ by $a^{-1}$ to obtain $b=0$.
\end{proof}

\begin{lemma}
For $0\leq a,b$ in an ordered field, $a\leq b$ if and only if $a^2\leq b^2$.
\end{lemma}
\begin{proof}
If $a\leq b$, then $a^2\leq ab\leq b^2$. Conversely, $a>b\geq0$ would give
$a^2-b^2=(a-b)(a+b)>0$, contradicting $a^2\leq b^2$.
\end{proof}

\begin{definition}[Absolute value]
For $x\in\mathbb R$,
\[
|x|=\begin{cases}x,&x\geq0,\\-x,&x<0.\end{cases}
\]
The quantity $|x-y|$ is the distance between $x$ and $y$.
\end{definition}
Useful identities include $|xy|=|x||y|$, $|x|^2=x^2$, and
$|x|<r\iff-r<x<r$ for $r>0$.

\begin{proposition}[Triangle and reverse triangle inequalities]
For $x,y\in\mathbb R$,
\[
|x+y|\leq|x|+|y|,\qquad
\bigl||x|-|y|\bigr|\leq|x-y|.
\]
\end{proposition}
\begin{proof}
Since $xy\leq|x||y|$, we have
$(x+y)^2\leq(|x|+|y|)^2$. Both sides to be compared after taking square
roots are nonnegative, giving the triangle inequality.
Applying it to $x=(x-y)+y$ gives $|x|-|y|\leq|x-y|$.
Interchanging $x,y$ gives the other half of the reverse inequality.
\end{proof}

\section{Completeness, supremum and infimum}
\begin{definition}[Bounds and extremal elements]
A set $S\subseteq\mathbb R$ is \emph{bounded above} if there is $M\in\mathbb R$
with $x\leq M$ for every $x\in S$. It is \emph{bounded below} if there is $m$
with $m\leq x$ for every $x\in S$. It is \emph{bounded} if both hold.

A maximum is an upper bound belonging to $S$; a minimum is a lower bound
belonging to $S$. Write $\max S$ and $\min S$ when they exist.
\end{definition}

\begin{definition}[Supremum and infimum]
For a nonempty set bounded above, $s=\sup S$ means:
\begin{enumerate}
\item $x\leq s$ for every $x\in S$;
\item for every $\varepsilon>0$, some $x\in S$ satisfies $s-\varepsilon<x$.
\end{enumerate}
Thus $s$ is the least upper bound. For a nonempty set bounded below,
$t=\inf S$ means that $t\leq x$ for all $x\in S$ and that every
$\varepsilon>0$ admits some $x\in S$ with $x<t+\varepsilon$.
\end{definition}

\begin{theorem}[Completeness of the real numbers]
There exists a complete ordered field, denoted $\mathbb R$.
Completeness means that every nonempty subset bounded above has a supremum
in $\mathbb R$. Equivalently, every nonempty subset bounded below has an
infimum in $\mathbb R$.
\end{theorem}
We take existence of this field as a foundational result; its construction
is outside these notes. The infimum statement follows from
$\inf S=-\sup(-S)$. Completeness is the extra property distinguishing
$\mathbb R$ from $\mathbb Q$.

\begin{proposition}
If $\max S$ exists, then $\sup S=\max S$. Similarly, if $\min S$ exists,
then $\inf S=\min S$.
\end{proposition}
\begin{proof}
A maximum $s\in S$ is an upper bound, and every upper bound is at least $s$
because it bounds every element of $S$. The minimum proof is analogous.
\end{proof}

\begin{theorem}[Archimedean property]\label{thm:analysis-archimedean}
The natural numbers are not bounded above in $\mathbb R$.
\end{theorem}
\begin{proof}
If $H=\sup\mathbb N$ existed, $H-1$ would not be an upper bound. Hence
some $n\in\mathbb N$ would satisfy $n>H-1$, and $n+1>H$ would contradict
that $H$ bounds $\mathbb N$.
\end{proof}
In particular, for every $\varepsilon>0$ there is $n\in\mathbb N$ with
$1/n<\varepsilon$.

\begin{example}[Existence of $\sqrt2$ from completeness]
Let $S=\{x\in\mathbb R:x>0,\ x^2<2\}$. It contains $1$ and is bounded
above by $2$, so $M=\sup S$ exists, with $1\leq M\leq2$.
We show $M^2=2$.

If $M^2<2$, choose $0<h<1$ so that
$(2M+1)h<2-M^2$. Then
\[
(M+h)^2=M^2+2Mh+h^2<M^2+(2M+1)h<2,
\]
contradicting that $M$ is an upper bound.

If $M^2>2$, choose $0<h<M$ with $2Mh<M^2-2$. Then
$(M-h)^2>M^2-2Mh>2$. Every $x\in S$ must satisfy $x<M-h$, since otherwise
$x^2\geq(M-h)^2>2$. Thus $M-h$ is an upper bound smaller than $M$, another
contradiction. Therefore $M^2=2$, and $M$ is the positive square root.
\end{example}

\begin{example}[Intervals and endpoints]
For $a<b$, the intervals are
\[
(a,b)=\{x:a<x<b\},\quad [a,b]=\{x:a\leq x\leq b\},\quad
[a,b)=\{x:a\leq x<b\}.
\]
\begin{center}
\begin{tabular}{lcccc}
\toprule
Set&Infimum&Minimum&Supremum&Maximum\\\midrule
$(a,b)$&$a$&none&$b$&none\\
$[a,b]$&$a$&$a$&$b$&$b$\\
$[a,b)$&$a$&$a$&$b$&none\\
$[a,\infty)$&$a$&$a$&none in $\mathbb R$&none\\\bottomrule
\end{tabular}
\end{center}
\end{example}

\begin{example}
For $S=\{3-1/n:n\in\mathbb N\}$, $\inf S=\min S=2$ and $\sup S=3$,
while no maximum exists. Indeed, all its elements are less than $3$.
Given $\varepsilon>0$, choose $n>1/\varepsilon$; then $3-1/n>3-\varepsilon$.
The same argument gives $\sup\{1-1/n:n\in\mathbb N\}=1$.
\end{example}

\begin{definition}[Neighbourhood]
For $a\in\mathbb R$ and $\varepsilon>0$, the $\varepsilon$-neighbourhood
of $a$ is $(a-\varepsilon,a+\varepsilon)=\{x:|x-a|<\varepsilon\}$.
\end{definition}
\begin{center}
\begin{tikzpicture}[>=Stealth]
\draw[->] (-3,0)--(3,0);
\draw[very thick,ink] (-1.8,0)--(1.8,0);
\foreach \x/\lab in {-1.8/{a-\varepsilon},0/a,1.8/{a+\varepsilon}}
 {\draw (\x,-.08)--(\x,.08) node[below=5pt] {$\lab$};}
\draw[fill=white] (-1.8,0) circle (2pt);
\draw[fill=white] (1.8,0) circle (2pt);
\end{tikzpicture}
\end{center}

\section{Sequences and their limits}
\begin{definition}[Sequence and boundedness]
A real sequence is a function $\mathbb N\to\mathbb R$, written $(x_n)_{n\geq1}$.
Its range is $\{x_n:n\in\mathbb N\}$. A sequence is bounded above, below,
or on both sides when its range has the corresponding property.
Equivalently, it is bounded if some $M>0$ satisfies $|x_n|\leq M$ for all $n$.
\end{definition}
For example, $x_n=n$ is bounded below but not above, whereas $x_n=(-1)^n$
is bounded. A finite initial part of a sequence does not determine its limit.

\begin{definition}[Convergence]\label{def:analysis-sequence-limit}
We say $x_n\to L\in\mathbb R$ if
\[
\forall\varepsilon>0\ \exists N\in\mathbb N\ \forall n\geq N:
\quad |x_n-L|<\varepsilon.
\]
We also write $\lim_{n\to\infty}x_n=L$. A sequence that has no finite real
limit is called divergent.
\end{definition}
The order of quantifiers matters: $N$ may depend on $\varepsilon$, but must
work for \emph{every} subsequent $n$. Replacing $n\geq N$ by $n>N$, or
allowing a real threshold $N$, gives an equivalent definition.

\begin{example}[Direct limit proofs]
For $x_n=1/n$, given $\varepsilon>0$, choose $N>1/\varepsilon$. Then
$n\geq N$ implies $|x_n|\leq1/N<\varepsilon$, so $1/n\to0$.

Also $1-2^{-n}\to1$, since $2^n\geq n$ and
$|(1-2^{-n})-1|=2^{-n}\leq1/n$.
For another example,
\[
\left|\frac{2n^2-1}{n^2+1}-2\right|=\frac3{n^2+1}\leq\frac3{n^2}.
\]
Choosing $N>\sqrt{3/\varepsilon}$ proves the limit is $2$.
\end{example}

\begin{remark}[Negating convergence to a specified value]
The statement $x_n\not\to L$ means
\[
\exists\varepsilon_0>0\ \forall N\in\mathbb N\ \exists n\geq N:
\quad |x_n-L|\geq\varepsilon_0.
\]
For example, $1/n\not\to1$: choose $\varepsilon_0=1/2$, and take
$n\geq\max\{N,2\}$. To prove divergence, one must exclude \emph{every}
possible finite limit, not merely one proposed value.
\end{remark}

\begin{theorem}[Uniqueness of the limit]
A sequence has at most one limit.
\end{theorem}
\begin{proof}
Suppose $x_n\to L_1$ and $x_n\to L_2$, with $L_1\ne L_2$. Set
$\varepsilon=|L_1-L_2|/2$. For sufficiently large $n$, both distances to
$x_n$ are less than $\varepsilon$. Thus
\[
|L_1-L_2|\leq|L_1-x_n|+|x_n-L_2|<2\varepsilon=|L_1-L_2|,
\]
a contradiction.
\end{proof}

\begin{theorem}[Convergent sequences are bounded]\label{thm:analysis-bounded}
If $x_n\to L\in\mathbb R$, then $(x_n)$ is bounded.
\end{theorem}
\begin{proof}
Use $\varepsilon=1$ to find $N$ with $|x_n-L|<1$ for $n\geq N$.
Then $|x_n|<|L|+1$ on this tail. The finitely many earlier terms are bounded
as well: take $M=\max\{|x_1|,\ldots,|x_N|,|L|+1\}$.
\end{proof}

\begin{theorem}[Sandwich theorem]\label{thm:analysis-sandwich}
Suppose $x_n\leq y_n\leq z_n$ for all sufficiently large $n$ and both
$x_n\to L$ and $z_n\to L$. Then $y_n\to L$.
\end{theorem}
\begin{proof}
For any $\varepsilon>0$, sufficiently large $n$ satisfy
$L-\varepsilon<x_n\leq y_n\leq z_n<L+\varepsilon$.
\end{proof}
\begin{example}
Since $-1/n^2\leq(-1)^n/n^2\leq1/n^2$, the sequence $(-1)^n/n^2$ tends
to $0$.
\end{example}

\section{Algebra of sequence limits}
\begin{theorem}[Algebra of limits]\label{thm:analysis-limit-algebra}
If $x_n\to x$ and $y_n\to y$, then
\[
x_n+y_n\to x+y,\qquad cx_n\to cx,\qquad x_ny_n\to xy.
\]
If additionally $y\ne0$, then $y_n\ne0$ eventually and $x_n/y_n\to x/y$
on that tail. If all $y_n\ne0$, the quotient is a sequence from the outset.
\end{theorem}
\begin{proof}
\textbf{Sum.} Make $|x_n-x|<\varepsilon/2$ and $|y_n-y|<\varepsilon/2$.
The triangle inequality gives
$|x_n+y_n-(x+y)|<\varepsilon$. The constant-multiple rule follows directly
from $|cx_n-cx|=|c||x_n-x|$ (and is immediate for $c=0$).

\textbf{Product.} By Theorem~\ref{thm:analysis-bounded}, choose $M\geq1$
with $|y_n|\leq M$ for all $n$. For sufficiently large $n$, both errors
$|x_n-x|$ and $|y_n-y|$ are less than $\varepsilon/(M+|x|)$. Then
\begin{align*}
|x_ny_n-xy|
&=|y_n(x_n-x)+x(y_n-y)|\\
&\leq M|x_n-x|+|x||y_n-y|<\varepsilon.
\end{align*}
\textbf{Quotient.} Since $y\ne0$, eventually
$|y_n-y|<|y|/2$, so $|y_n|>|y|/2$. Hence
\[
\left|\frac1{y_n}-\frac1y\right|
=\frac{|y_n-y|}{|y_n||y|}
\leq\frac{2|y_n-y|}{|y|^2}\longrightarrow0.
\]
Apply the product rule to $x_n(1/y_n)$.
\end{proof}
\begin{remark}
If $y_n\ne0$ for every $n$ and $y_n\to y\ne0$, then in fact $|y_n|\geq c>0$
for every $n$. Take the minimum of $|y|/2$ and the finitely many nonzero
initial absolute values. This is the source's auxiliary quotient-rule lemma.
\end{remark}
\begin{example}
Dividing numerator and denominator by $n^2$ gives
\[
\frac{2n^2-1}{n^2+1}=\frac{2-1/n^2}{1+1/n^2}\longrightarrow2.
\]
This shorter argument becomes valid once the algebra of limits is proved.
\end{example}
\begin{proposition}[Order passes to limits]
If $x_n\leq y_n$ eventually, $x_n\to x$, and $y_n\to y$, then $x\leq y$.
\end{proposition}
\begin{proof}
If $x>y$, choose $\varepsilon=(x-y)/3$. Eventually
$x_n>x-\varepsilon>y+\varepsilon>y_n$, a contradiction.
\end{proof}

\section{Infinite limits and monotone sequences}
\begin{definition}[Divergence to infinity]
We write $x_n\to+\infty$ if for every $M\in\mathbb R$ some $N$ satisfies
$x_n>M$ whenever $n\geq N$. We write $x_n\to-\infty$ if for every
$m\in\mathbb R$ eventually $x_n<m$.
\end{definition}
These are forms of divergence, not convergence to a real number. For example,
$n^2\to+\infty$, but $(-1)^n$ has neither a finite nor an infinite limit.

\begin{definition}[Monotonicity]
A sequence is \emph{increasing} (nondecreasing) if $x_{n+1}\geq x_n$ for
all $n$, and \emph{strictly increasing} if these inequalities are strict.
Decreasing and strictly decreasing are defined by reversing the signs.
A sequence is \emph{monotone} if it is increasing or decreasing.
\end{definition}

\begin{theorem}[Monotone convergence]\label{thm:analysis-monotone}
An increasing sequence bounded above converges to $\sup\{x_n:n\in\mathbb N\}$.
A decreasing sequence bounded below converges to its infimum.
\end{theorem}
\begin{proof}
In the increasing case, put $s=\sup\{x_n\}$. Given $\varepsilon>0$, some
$x_N>s-\varepsilon$. For every $n\geq N$,
$s-\varepsilon<x_N\leq x_n\leq s<s+\varepsilon$.
Apply this result to $(-x_n)$ to obtain the decreasing case.
\end{proof}
An increasing sequence unbounded above tends to $+\infty$: once it exceeds
any prescribed $M$, every subsequent term does too. The decreasing analogue
holds for $-\infty$.

\begin{example}[A recursive limit]
Let $x_1=1$ and $x_{n+1}=(x_n+1)/3$. Induction gives $x_n\geq1/2$, because
$x_n\geq1/2$ implies $x_{n+1}\geq(1/2+1)/3=1/2$. Also
\[
x_{n+1}-x_n=\frac23\left(\frac12-x_n\right)\leq0.
\]
The sequence is decreasing and bounded below, so it converges to some $L$.
Passing to the limit gives $L=(L+1)/3$, hence $L=1/2$.
\end{example}
\begin{remark}[Prove existence before solving for the limit]
For the Fibonacci sequence $F_1=F_2=1$ and $F_{n+2}=F_{n+1}+F_n$, blindly
assuming a finite limit gives $L=2L$, hence $L=0$. This is invalid: the
sequence is unbounded. Indeed $F_{n+2}\geq F_{n+1}+1$, so it grows without
bound. A fixed-point equation identifies a possible limit only after
convergence has been justified.
\end{remark}

\section{Subsequences and Bolzano--Weierstrass}
\begin{definition}[Subsequence]
If $j_1<j_2<\cdots$ are natural numbers, then $(x_{j_n})$ is a subsequence
of $(x_n)$. Necessarily $j_n\geq n$.
\end{definition}
\begin{theorem}
If $x_n\to L$, every subsequence tends to $L$. The same holds for
$L=+\infty$ or $-\infty$.
\end{theorem}
\begin{proof}
Once all terms with index at least $N$ satisfy the required estimate,
$x_{j_n}$ satisfies it for $n\geq N$, since $j_n\geq n$.
\end{proof}
Thus two subsequences with different limits prove divergence. In particular,
a fixed shift $x_{n+k}$ has the same limit as $x_n$.

\begin{example}
For $x_n=1/n+\sin(n\pi/2)$, we have $x_{2n}=1/(2n)\to0$ while
$x_{4n+1}=1/(4n+1)+1\to1$. Therefore $(x_n)$ diverges.
\end{example}

\begin{lemma}
Every real sequence has a monotone subsequence.
\end{lemma}
\begin{proof}
Call $n$ a \emph{peak} if $x_n\geq x_m$ for every $m>n$.
If there are infinitely many peaks, list them in increasing order
$j_1<j_2<\cdots$. Then $(x_{j_n})$ is decreasing.

If there are only finitely many peaks, choose $j_1$ beyond all of them
(or $j_1=1$ if there are none). As $j_1$ is not a peak, some $j_2>j_1$
satisfies $x_{j_2}>x_{j_1}$. Repeating constructs a strictly increasing
subsequence.
\end{proof}

\begin{theorem}[Bolzano--Weierstrass]\label{thm:analysis-bw}
Every bounded real sequence has a convergent subsequence.
\end{theorem}
\begin{proof}
Take a monotone subsequence using the lemma. It is still bounded, so it
converges by Theorem~\ref{thm:analysis-monotone}.
\end{proof}
\begin{example}[Many subsequential limits]
List the fractions in $(0,1)$ in successive denominator blocks:
\[
\frac12;\quad \frac13,\frac23;\quad
\frac14,\frac24,\frac34;\quad\ldots.
\]
Every $r\in[0,1]$ is the limit of a subsequence. For each denominator
$q\geq2$, choose $p_q\in\{1,\ldots,q-1\}$ with
$|p_q/q-r|\leq1/q$. Picking this entry from each consecutive block gives
a subsequence tending to $r$. This supplies the exercise argument on
source p.~32.
\end{example}

\section{The Cauchy--Schwarz inequality}
\begin{theorem}[Cauchy--Schwarz]
For real $x_1,\ldots,x_n$ and $y_1,\ldots,y_n$,
\[
\left(\sum_{k=1}^n x_ky_k\right)^2
\leq\left(\sum_{k=1}^n x_k^2\right)
     \left(\sum_{k=1}^n y_k^2\right).
\]
\end{theorem}
\begin{proof}
Write $A=\sum y_k^2$, $B=\sum x_ky_k$ and $C=\sum x_k^2$.
If $A=0$, every $y_k=0$ and the claim is immediate. Otherwise,
\[
0\leq\sum_{k=1}^n(x_k-ty_k)^2=At^2-2Bt+C
\]
for every real $t$. Substituting $t=B/A$ gives $C-B^2/A\geq0$, proving
$B^2\leq AC$.
\end{proof}
For $n=2$, the difference between the right and left sides is
$(x_1y_2-x_2y_1)^2$. In $\mathbb R^3$, the geometric interpretation is
$|\mathbf x\cdot\mathbf y|\leq\|\mathbf x\|\|\mathbf y\|$, since the dot
product equals $\|\mathbf x\|\|\mathbf y\|\cos\theta$.
